% BEGIN VOL07-EXPANSION VII32
\section{Supplementary worked examples}
\label{sec:vii32-pedagogy}

\begin{example}[Vertical logarithmic distance]\label{ex:vii32-ped-01}
Along a vertical geodesic, \(d(iy_1,iy_2)=|\log(y_2/y_1)|\).  Equal multiplicative changes in height therefore correspond to equal additive hyperbolic distances.
\end{example}

\begin{example}[Cusp strip area]\label{ex:vii32-ped-02}
The hyperbolic area of \(0\le x\le1\), \(y\ge Y\) is \(\int_Y^\infty\int_0^1 y^{-2}\,dx\,dy=1/Y\).  A region extending infinitely high can have finite hyperbolic area.
\end{example}

\begin{example}[Boundary at infinite distance]\label{ex:vii32-ped-03}
The vertical distance from \(i\) to \(iy\) is \(|\log y|\), which diverges as \(y\downarrow0\) or \(y\to\infty\).  Euclidean boundary points are therefore infinitely far away intrinsically.
\end{example}

\section{Graded supplementary exercises}
The following exercises add a deliberate progression from concrete calculation to structural proof, hypothesis testing, application, and synthesis.

\subsection*{Standard computations and constructions}

\begin{exercise}[Vertical distance]\label{exr:vii32-ped-01}
Compute \(d(i,2i)\).
\end{exercise}

\begin{hint}
Use the logarithmic formula.
\end{hint}

\begin{solution}
\(\log2\).
\end{solution}

\begin{exercise}[Another vertical distance]\label{exr:vii32-ped-02}
Compute \(d(2i,8i)\).
\end{exercise}

\begin{hint}
Take the height ratio.
\end{hint}

\begin{solution}
\(\log4\).
\end{solution}

\begin{exercise}[Area density]\label{exr:vii32-ped-03}
What is the area form of the half-plane metric?
\end{exercise}

\begin{hint}
Take the determinant of \(y^{-2}I\).
\end{hint}

\begin{solution}
\(dx\,dy/y^2\).
\end{solution}

\begin{exercise}[Curve speed]\label{exr:vii32-ped-04}
For \(\gamma(t)=i e^t\), compute hyperbolic speed.
\end{exercise}

\begin{hint}
\(|\dot z|/\operatorname{Im}z\).
\end{hint}

\begin{solution}
The speed is \(1\).
\end{solution}

\begin{exercise}[Distance formula]\label{exr:vii32-ped-05}
Compute \(\cosh d(i,1+i)\).
\end{exercise}

\begin{hint}
Use the general formula.
\end{hint}

\begin{solution}
\(1+1/2=3/2\).
\end{solution}

\subsection*{Proofs}

\begin{exercise}[Boundary distance]\label{exr:vii32-ped-06}
Prove the boundary \(y=0\) is at infinite vertical distance.
\end{exercise}

\begin{hint}
Integrate \(dy/y\) to zero.
\end{hint}

\begin{solution}
The integral diverges logarithmically.
\end{solution}

\begin{exercise}[Conformal angles]\label{exr:vii32-ped-07}
Why do Euclidean and hyperbolic angles agree in the half-plane model?
\end{exercise}

\begin{hint}
The metric is a scalar multiple of Euclidean metric.
\end{hint}

\begin{solution}
Conformal scaling preserves angle cosines.
\end{solution}

\begin{exercise}[Cusp area]\label{exr:vii32-ped-08}
Derive the area \(1/Y\) of the strip \([0,1]\times[Y,\infty)\).
\end{exercise}

\begin{hint}
Integrate \(y^{-2}\).
\end{hint}

\begin{solution}
\(\int_Y^\infty y^{-2}dy=1/Y\).
\end{solution}

\begin{exercise}[Scaling isometry]\label{exr:vii32-ped-09}
Show \(z\mapsto\lambda z\), \(\lambda>0\), preserves the Poincare line element.
\end{exercise}

\begin{hint}
Both numerator and height scale by \(\lambda\).
\end{hint}

\begin{solution}
\(|d(\lambda z)|/\operatorname{Im}(\lambda z)=|dz|/\operatorname{Im}z\).
\end{solution}

\subsection*{Counterexamples and hypothesis tests}

\begin{exercise}[Horizontal lines geodesic]\label{exr:vii32-ped-10}
Test the claim in the upper half-plane.
\end{exercise}

\begin{hint}
Known geodesics meet the real boundary orthogonally.
\end{hint}

\begin{solution}
False.  Horizontal Euclidean lines are not hyperbolic geodesics.
\end{solution}

\begin{exercise}[Raw Euclidean distance]\label{exr:vii32-ped-11}
Test the claim: Euclidean distance between coordinates equals hyperbolic distance.
\end{exercise}

\begin{hint}
Compare \(i\) and \(2i\).
\end{hint}

\begin{solution}
False.  Their Euclidean distance is \(1\), while hyperbolic distance is \(\log2\).
\end{solution}

\begin{exercise}[Boundary finite]\label{exr:vii32-ped-12}
Test the claim: points arbitrarily close to \(y=0\) are at bounded distance from \(i\).
\end{exercise}

\begin{hint}
Use vertical distance.
\end{hint}

\begin{solution}
False.  The distance tends to infinity as height tends to zero.
\end{solution}

\subsection*{Applications and investigations}

\begin{exercise}[Cusp finite area]\label{exr:vii32-ped-13}
Why can a hyperbolic cusp have finite area but infinite extent?
\end{exercise}

\begin{hint}
Compare the decay of the area density with the distance integral.
\end{hint}

\begin{solution}
Area integrates \(y^{-2}\), which is integrable at infinity, while vertical distance integrates \(y^{-1}\), which diverges.
\end{solution}

\begin{exercise}[Geodesic drawing]\label{exr:vii32-ped-14}
How should a half-plane geodesic through two generic points be drawn?
\end{exercise}

\begin{hint}
Use the geodesic-shape theorem.
\end{hint}

\begin{solution}
As the unique Euclidean circle through them whose center lies on the real axis, unless they share the same real part, in which case it is vertical.
\end{solution}

\subsection*{Challenge problems}

\begin{exercise}[Unit-speed vertical]\label{exr:vii32-ped-15}
Show \(\gamma(t)=ie^t\) is unit speed and travels distance \(|t_2-t_1|\).
\end{exercise}

\begin{hint}
Compute speed and integrate.
\end{hint}

\begin{solution}
Its speed is \(1\), so length equals the parameter difference.
\end{solution}

\begin{exercise}[Distance inversion]\label{exr:vii32-ped-16}
If \(d(iy_1,iy_2)=L\) with \(y_2>y_1\), express \(y_2/y_1\).
\end{exercise}

\begin{hint}
Exponentiate the logarithmic distance.
\end{hint}

\begin{solution}
\(y_2/y_1=e^L\).
\end{solution}

% END VOL07-EXPANSION VII32
